\chapter{Literature Survey}

\section{Introduction}
This chapter reviews existing literature on Virtual Reality (VR) in education and training, specifically focusing on aviation safety.

\section{VR in Safety Training}
Research has shown that VR is highly effective for safety training. Chittaro et al. (2014) demonstrated that interactive VR safety briefings significantly improved knowledge retention compared to traditional safety cards. The immersive nature of VR allows users to build "muscle memory" for procedures that they would otherwise only read about.

\section{Limitations of Current Methods}
Traditional pre-flight briefings are often ignored. A study by the NTSB (National Transportation Safety Board) highlighted that many passengers involved in accidents were unable to recall basic safety instructions. This suggests a need for more engaging training methods.

\section{StereoKit Framework}
StereoKit is a lightweight, open-source Mixed Reality library for C\#. Unlike heavy game engines like Unity or Unreal, StereoKit is designed for building mixed reality tools and applications with a code-first approach. It provides a high-performance rendering pipeline and a simplified API for handling VR inputs, making it an ideal choice for this project.

\section{Conclusion}
The literature supports the hypothesis that VR can improve safety training outcomes. However, most existing solutions are built on heavy game engines. This project explores the use of a lightweight framework (StereoKit) to deliver a high-fidelity training experience on standalone hardware.
